\section*{Capítulo \ref{ch:g10:lewis-and-shape} --- Estruturas de Lewis e geometria das moléculas}

\begin{solution}{exo:g10:lewis-and-shape:1}
Um par de \omterm{def:g9:inside-the-atom:nucleus}{elétrons} compartilhado por dois \omterm{def:g7:atoms-and-molecules:atom}{átomos}. Um par de \omterm{def:g10:electron-shells:valence}{elétrons de
valência} que pertence a um só \omterm{def:g7:atoms-and-molecules:atom}{átomo}, sem ser compartilhado.
\end{solution}

\begin{solution}{exo:g10:lewis-and-shape:2}
H: 1 (falta-lhe um \omterm{def:g9:inside-the-atom:nucleus}{elétron} para o dueto). C: 4, N: 3, O: 2 (faltam-lhes
4, 3 e 2 \omterm{def:g9:inside-the-atom:nucleus}{elétrons} para o octeto).
\end{solution}

\begin{solution}{exo:g10:lewis-and-shape:3}
\chemfig{H-\charge{0=\:,90=\:,270=\:}{Cl}}: o cloro tem o par compartilhado e três \omterm{def:g10:lewis-and-shape:lone-pair}{pares
não ligantes}, 8 \omterm{def:g9:inside-the-atom:nucleus}{elétrons}.
\end{solution}

\begin{solution}{exo:g10:lewis-and-shape:4}
$5 + 3 \times 1 = 8$ \omterm{def:g9:inside-the-atom:nucleus}{elétrons}, 4 pares (três ligações, um \omterm{def:g10:lewis-and-shape:lone-pair}{par não ligante}).
\end{solution}

\begin{solution}{exo:g10:lewis-and-shape:5}
\ce{CH4}: tetraédrica. \ce{NH3}: piramidal. \ce{H2O}: angular.
\ce{CO2}: linear.
\end{solution}

\begin{solution}{exo:g10:lewis-and-shape:6}
\chemfig{H-\charge{90=\:,270=\:}{S}-H}: como na água, quatro grupos em volta do enxofre,
dos quais dois \omterm{def:g10:lewis-and-shape:lone-pair}{pares não ligantes}: angular.
\end{solution}

\begin{solution}{exo:g10:lewis-and-shape:7}
$2 \times 5 = 10$ \omterm{def:g9:inside-the-atom:nucleus}{elétrons}, 5 pares. Uma ligação simples e 4 \omterm{def:g10:lewis-and-shape:lone-pair}{pares não
ligantes} deixam cada nitrogênio com 6 \omterm{def:g9:inside-the-atom:nucleus}{elétrons}; transformar dois \omterm{def:g10:lewis-and-shape:lone-pair}{pares
não ligantes} em pares compartilhados dá uma \omterm{def:g10:lewis-and-shape:multiple-bond}{ligação tripla} e um \omterm{def:g10:lewis-and-shape:lone-pair}{par não
ligante} em cada \omterm{def:g7:atoms-and-molecules:atom}{átomo}:
\chemfig{\charge{180=\:}{N}~\charge{0=\:}{N}}.
\end{solution}

\begin{solution}{exo:g10:lewis-and-shape:8}
\chemfig{H-C(-[2]H)(-[6]H)-\charge{90=\:,270=\:}{O}-H}: quatro ligações em volta do
carbono, tetraédrica; em volta do oxigênio, duas ligações e dois \omterm{def:g10:lewis-and-shape:lone-pair}{pares
não ligantes}, angular.
\end{solution}

\begin{solution}{exo:g10:lewis-and-shape:9}
No \ce{CO2}, o carbono tem dois grupos (duas \omterm{def:g10:lewis-and-shape:multiple-bond}{ligações duplas}) e nenhum
\omterm{def:g10:lewis-and-shape:lone-pair}{par não ligante}: eles apontam em sentidos opostos. No \ce{H2O}, o
oxigênio tem quatro grupos (duas ligações, dois \omterm{def:g10:lewis-and-shape:lone-pair}{pares não ligantes}) numa
disposição tetraédrica: as duas ligações formam um ângulo.
\end{solution}

\begin{solution}{exo:g10:lewis-and-shape:10}
\chemfig{C(-[3]H)(-[5]H)=C(-[1]H)-[7]H}: em volta de cada carbono, três
grupos (uma \omterm{def:g10:lewis-and-shape:multiple-bond}{ligação dupla}, duas ligações simples): trigonal plana,
ângulos perto de \ang{120}.
\end{solution}

\begin{solution}{exo:g10:lewis-and-shape:11}
O carbono no centro, quatro ligações simples com \omterm{def:g7:atoms-and-molecules:atom}{átomos} de cloro, cada
um com três \omterm{def:g10:lewis-and-shape:lone-pair}{pares não ligantes}: quatro grupos em volta do carbono,
geometria tetraédrica.
\end{solution}

\begin{solution}{exo:g10:lewis-and-shape:12}
O etanol, \chemfig{H-C(-[2]H)(-[6]H)-C(-[2]H)(-[6]H)-\charge{90=\:,270=\:}{O}-H}, e o
metoximetano,
\chemfig{H-C(-[2]H)(-[6]H)-\charge{90=\:,270=\:}{O}-C(-[2]H)(-[6]H)-H}.
\end{solution}

\begin{solution}{exo:g10:lewis-and-shape:13}
Um \omterm{def:g10:lewis-and-shape:lone-pair}{par não ligante} ocupa um pouco mais de espaço que um \omterm{def:g10:lewis-and-shape:lone-pair}{par ligante} e
aproxima as ligações. A amônia tem um \omterm{def:g10:lewis-and-shape:lone-pair}{par não ligante}, a água dois: o
ângulo diminui do metano para a amônia e da amônia para a água.
\end{solution}

\begin{solution}{exo:g10:lewis-and-shape:14}
\chemfig{H-C~\charge{0=\:}{N}}: dois grupos em volta do carbono, linear.
\chemfig{H-C(-[6]H)=\charge{45=\:,315=\:}{O}}: três grupos em volta do carbono,
trigonal plana.
\end{solution}

\begin{solution}{exo:g10:lewis-and-shape:15}
$3 \times 6 = 18$ \omterm{def:g9:inside-the-atom:nucleus}{elétrons}, 9 pares:
\chemfig{\charge{180=\:,270=\:}{O}=\charge{90=\:}{O}-\charge{0=\:,90=\:,270=\:}{O}}. O oxigênio
central tem três grupos (uma \omterm{def:g10:lewis-and-shape:multiple-bond}{ligação dupla}, uma ligação simples, um \omterm{def:g10:lewis-and-shape:lone-pair}{par
não ligante}): a \omterm{def:g7:atoms-and-molecules:molecule}{molécula} é angular.
\end{solution}

\begin{solution}{pb:g10:lewis-and-shape:1}
\textbf{1.} \chemfig{H-C(-[2]H)(-[6]H)-H}, \chemfig{H-\charge{90=\:}{N}(-[6]H)-H},
\chemfig{H-\charge{90=\:,270=\:}{O}-H}.

\textbf{2.} Quatro grupos em cada uma.

\textbf{3.} Nenhum no metano, um na amônia, dois na água.

\textbf{4.} O carbono forma quatro ligações; o oxigênio forma duas (os
seus dois outros grupos são \omterm{def:g10:lewis-and-shape:lone-pair}{pares não ligantes}, que um kit não mostra
com varetas).

\textbf{5.} \chemfig{H-C(-[2]H)(-[6]H)-C(-[2]H)(-[6]H)-H}: uma vareta liga
as duas bolas de carbono, e são necessárias seis bolas de hidrogênio.

\textbf{6.} Tetraédrica, piramidal, angular.

\textbf{7.} Os \omterm{def:g10:lewis-and-shape:lone-pair}{pares não ligantes} ocupam mais espaço que os \omterm{def:g10:lewis-and-shape:lone-pair}{pares
ligantes} e apertam as ligações umas contra as outras; a água, com dois
\omterm{def:g10:lewis-and-shape:lone-pair}{pares não ligantes}, tem o menor ângulo.

\textbf{8.} Os furos de uma bola de carbono formam \ang{109.5}; o ângulo
verdadeiro da água é \ang{104.5}: uma bola de oxigênio precisa de seus
próprios dois furos a \ang{104.5}.

\textbf{9.} \ang{180}: dois grupos em volta do carbono apontam em
sentidos opostos, e a \omterm{def:g7:atoms-and-molecules:molecule}{molécula} é linear.

\textbf{10.} Dois vértices que não estão na mesma aresta são vértices
opostos de uma face: a distância entre eles é a diagonal da face,
$\sqrt{a^2 + a^2} = a\sqrt{2}$.

\textbf{11.} O centro fica no meio de uma diagonal do cubo:
$a\sqrt{3}/2$.

\textbf{12.} H--H: $\sqrt{2} \approx \qty{1.41}{cm}$; C--H:
$\sqrt{3}/2 \approx \qty{0.87}{cm}$.

\textbf{13.} Cada hidrogênio está num vértice do cubo e o carbono no seu
centro: todos os vértices estão à mesma distância do centro.

\textbf{14.} O triângulo C\,H$_1$\,H$_2$ é isósceles (C\,H$_1$ = C\,H$_2$);
a reta que vai do vértice ao ponto médio da base é perpendicular à base.

\textbf{15.} $\sin \widehat{\text{M\,C\,H}_1} = \dfrac{\text{M\,H}_1}{\text{C\,H}_1}
= \dfrac{a\sqrt{2}/2}{a\sqrt{3}/2} = \sqrt{2/3} \approx 0.816$.

\textbf{16.} Meio ângulo $\approx \ang{54.7}$; ângulo H--C--H $\approx
2 \times 54.7 = \ang{109.5}$.

\textbf{17.} É o ângulo de \ang{109.5} do tetraedro da tabela.

\textbf{18.} Os furos devem ser feitos a \ang{109.5} uns dos outros.
\end{solution}
